\chapter{Signaux tirés des options pour les actions}\label{ch:s1:options-implied-signals-for-stocks}

Lorsque les puts hors de la monnaie d'une action deviennent chers par rapport à ses options à la monnaie, quelqu'un paie pour se protéger, et l'action tend à baisser.
Xing, Zhang et Zhao ont montré que les actions aux smirks de volatilité les plus pentus ont sous-performé celles aux plus plats de 10.9~\% par an, ajustés du risque, et ont eu
les pires surprises de résultats le trimestre suivant~: les traders informés porteurs de mauvaises nouvelles préfèrent les puts. Le marché des options est celui où le levier est bon marché et où la vente à découvert est
superflue, c'est donc là qu'une partie de l'information apparaît en premier. Ce chapitre ajoute une couche d'options au marché synthétique, avec des traders informés qui achètent des options
avant des annonces de résultats programmées, et mesure ce que prédisent le skew, l'écart de volatilité call--put et le volume d'options, à quelle vitesse ce signal décroît, et pourquoi il est préférable de le négocier
autour des événements. La construction est \texttt{firm.optsignal}.

\section{Le skew comme signal}

\begin{definition}[Trading d'options informé]\label{def:s1:options-implied-signals-for-stocks:informed}
Le \emph{trading d'options informé}\index{trading d'options informé} désigne l'activité sur les options d'une action par des investisseurs qui détiennent une information privée sur celle-ci, et qui
préfèrent les options à l'action pour leur effet de levier et, en cas de mauvaise nouvelle, pour éviter les coûts et les contraintes de la vente à découvert~; leur demande fait bouger les prix et les
volumes des options qu'ils achètent avant que l'information n'atteigne l'action.
\end{definition}

Le skew de volatilité (livre~1, chapitre~25) est la volatilité implicite supplémentaire des puts hors de la monnaie par rapport aux options à la monnaie. Il existe sur chaque action (une
prime de protection contre les krachs), et il se déplace lorsque les acheteurs de puts sont plus nombreux que les vendeurs. Dans la couche d'options synthétique, chaque action a une volatilité implicite à la monnaie
égale à sa volatilité réalisée, augmentée d'une prime de deux points, d'un skew de put de trois points et d'un bruit quotidien d'un point sur chacune~; les traders informés,
dans les dix jours qui précèdent chaque annonce de résultats (calendrier du chapitre~7), achètent des puts avant les mauvaises surprises et des calls avant les bonnes, et chaque
écart-type de surprise déplace la volatilité implicite des options qu'ils achètent de 0.3 point. Les mauvaises nouvelles sont pondérées 1.5 fois plus que les bonnes, pour la raison
que Johnson et So exposent plus bas. Le signe du skew comme signal est négatif~: un skew pentu annonce une baisse.

\section{Volatilité implicite moins volatilité réalisée}

\begin{definition}[Écart de volatilité implicite]\label{def:s1:options-implied-signals-for-stocks:spread}
Le \emph{écart de volatilité implicite}\index{écart de volatilité implicite} d'une action est la différence entre les volatilités implicites de ses calls et de ses puts à
strikes et échéances égaux~; sous la parité call--put, il est nul, et un écart positif signifie que les calls sont relativement chers.
\end{definition}

Cremers et Weinbaum ont mesuré les écarts à la parité call--put de cette façon et ont constaté que les actions aux calls relativement chers ont surperformé celles
dont les puts étaient relativement chers de 50 points de base par semaine, un écart non expliqué par les contraintes de vente à découvert, plus marqué là où les options étaient liquides et les actions ne l'étaient pas, et
décroissant sur leur échantillon. Bali et Hovakimian ont séparé deux écarts~: l'écart call--put (positivement lié aux rendements futurs, un indicateur du risque de saut et du
trading informé) et l'écart implicite-réalisé (négativement lié aux rendements par le biais de réalisé moins implicite, un indicateur du risque de volatilité)~;
les niveaux de volatilité ne prédisaient rien. Dans la couche synthétique, l'écart implicite-réalisé ne porte aucune information injectée, et son IC est nul à tous les horizons~:
un contrôle utile.

\begin{center}
\small
\setlength{\tabcolsep}{4pt}%
\begin{tabular}{@{}lcccc@{}}
\toprule
IC de rang quotidien moyen, années 3 à 10 & 1 jour & 5 jours & 20 jours & 60 jours\\
\midrule
skew (signe moins), toutes les actions & 0.004 & 0.014 & 0.018 & 0.012\\
écart de volatilité call--put, toutes les actions & 0.005 & 0.017 & 0.022 & 0.015\\
\omterm{def:s1:options-implied-signals-for-stocks:volume}{ratio de volume options sur actions} (signe moins), toutes les actions & 0.002 & 0.008 & 0.010 & 0.007\\
volatilité implicite moins réalisée, toutes les actions & 0.000 & 0.001 & 0.000 & $-0.000$\\
\midrule
skew (signe moins), annonçant dans les dix jours & 0.026 & 0.085 & 0.111 & 0.082\\
écart de volatilité call--put, annonçant dans les dix jours & 0.030 & 0.101 & 0.129 & 0.094\\
volume options sur actions (signe moins), annonçant dans les dix jours & 0.015 & 0.047 & 0.061 & 0.045\\
\bottomrule
\end{tabular}
\end{center}

\section{Volume d'options et ratio put--call}

\begin{definition}[Ratio put--call, ratio de volume options sur actions]\label{def:s1:options-implied-signals-for-stocks:volume}
Le \emph{ratio put--call}\index{ratio put--call} d'une action est le volume de ses options de vente divisé par le volume de ses options d'achat sur une période, idéalement en ne comptant que
le volume qui ouvre de nouvelles positions longues. Son \emph{ratio de volume options sur actions}\index{ratio de volume options sur actions} est son volume total d'options (en nombre d'actions
sous-jacentes) divisé par son volume d'actions.
\end{definition}

Pan et Poteshman ont construit des \omterm{def:s1:options-implied-signals-for-stocks:volume}{ratios put--call} à partir du volume initié par les acheteurs qui ouvraient de nouvelles positions, une base de données dont la plupart des chercheurs ne disposent pas, et ont constaté que
les actions aux ratios faibles devancent celles aux ratios élevés de plus de 40 points de base le lendemain et de plus de 1~\% sur la semaine suivante, grâce à une information non
publique détenue par les traders d'options. Johnson et So ont montré pourquoi le volume non signé est aussi informatif~: les coûts de vente à découvert poussent les traders informés à utiliser davantage les options
pour les mauvaises nouvelles que pour les bonnes, si bien qu'un \omterm{def:s1:options-implied-signals-for-stocks:volume}{ratio de volume options sur actions} élevé signale une mauvaise nouvelle~; le décile le plus bas a devancé le plus haut de 0.34~\% par semaine. Dans la couche
synthétique, le volume d'options croît avec la taille de la demande informée quelle que soit sa direction, avec une pondération vers les mauvaises nouvelles, si bien que son IC vaut environ la moitié de celui des signaux de
prix~: le volume indique que quelqu'un sait quelque chose, et l'asymétrie dit de quoi il s'agit.

\section{Où les traders informés opèrent}

La moitié inférieure du tableau est la leçon du chapitre. Sur l'ensemble des actions, les IC des signaux sont faibles (0.02 pour l'écart de volatilité à vingt jours), car la plupart
des actions, à la plupart des dates, n'ont pas de traders informés dans leurs options, et les signaux ne sont que du bruit. Parmi les actions annonçant dans les dix jours (11.7~\% des jours-actions), les
mêmes signaux ont des IC environ six fois plus élevés~: l'information est concentrée là où se trouvent les événements. Les IC montent jusqu'à vingt jours puis décroissent, puisque le
saut de l'annonce, que les traders informés avaient anticipé, tombe dans la fenêtre de vingt jours, et la dérive post-annonce injectée en emporte une partie
plus loin.

\begin{omfigure}
\begin{tikzpicture}
\begin{axis}[width=10.5cm,height=5.4cm,xmode=log,xlabel={\small horizon (jours)},ylabel={\small IC de rang quotidien moyen},tick label style={font=\footnotesize},
  xtick={1,5,20,60},xticklabels={1,5,20,60},xmin=0.8,xmax=75,ymin=0,ymax=0.14,ytick={0,0.05,0.1},
  yticklabel style={/pgf/number format/fixed,/pgf/number format/precision=2},grid=major,grid style={gray!25},
  legend style={font=\scriptsize,at={(1.02,1)},anchor=north west},table/col sep=comma]
\addplot[omDef,very thick,mark=*] table[x=h,y=spread_ev]{figdata/strategies-1/15-options-implied-signals-for-stocks/ic.csv};
\addplot[omThm,very thick,mark=square*] table[x=h,y=skew_ev]{figdata/strategies-1/15-options-implied-signals-for-stocks/ic.csv};
\addplot[black!50,thick,mark=triangle*] table[x=h,y=os_ev]{figdata/strategies-1/15-options-implied-signals-for-stocks/ic.csv};
\addplot[omDef,thick,dashed,mark=o] table[x=h,y=spread]{figdata/strategies-1/15-options-implied-signals-for-stocks/ic.csv};
\addplot[omThm,thick,dashed,mark=square] table[x=h,y=skew]{figdata/strategies-1/15-options-implied-signals-for-stocks/ic.csv};
\legend{écart (annonces),skew (annonces),volume (annonces),écart (ensemble),skew (ensemble)}
\end{axis}
\end{tikzpicture}

\omcaption{IC de rang des signaux d'options contre le rendement ajusté du marché sur les 1 à 60 jours suivants, sur le marché synthétique avec des traders
d'options informés~: parmi les actions annonçant dans les dix jours (trait plein) et toutes les actions (tirets). Données~: \texttt{s1\_options.ic}.}
\label{fig:s1:options-implied-signals-for-stocks:ic}
\end{omfigure}

Un portefeuille hebdomadaire par déciles sur l'écart de volatilité, sur l'ensemble des actions, rapporte 8.5~\% par an avant coûts (ratio de Sharpe 3.67) et perd 0.9~\% après dix
points de base par unité échangée~: le bruit quotidien des volatilités implicites fait tourner le classement. Restreint aux actions annonçant dans les dix jours (quintiles de ce
sous-ensemble), il rapporte 31.9~\% par an après coûts, avec un ratio de Sharpe de 6.12, le portefeuille de skew 23.0~\% et le portefeuille de volume 9.1~\%. Ces chiffres synthétiques sont très
au-dessus de tout ce que la littérature publie, et la raison tient au modèle~: ses traders informés connaissent toute la surprise et leur demande en est une fonction nette. Les flux informés réels
sont une part faible et non signée du volume d'options, partagée avec les hedgers et les spéculateurs~; ce qui se transpose, c'est la structure~: négocier les signaux là où
se trouvent les événements, et rééquilibrer seulement aussi vite que l'information change.

\section{Fiches de stratégie}

\begin{strategyfile}{Signal de skew}\label{strat:s1:options-implied-signals-for-stocks:skew}
\sfield{Qui paie, et pourquoi} Traders informés qui révèlent de mauvaises nouvelles en achetant des puts, et la lenteur du marché des actions à suivre.
\sfield{Instruments et plateformes} Actions dont les options sont liquides~; le signal provient des options, la transaction se fait sur l'action.
\sfield{Signal} Volatilité implicite des puts hors de la monnaie moins celle des options à la monnaie, rapportée à l'historique propre de l'action.
\sfield{Taille et exécution} Vendre les skews pentus, acheter les plats~; concentrer les positions autour des événements programmés.
\sfield{Coûts} Transactions sur actions~; le bruit du signal fait monter la rotation du portefeuille.
\sfield{Comment elle s'éteint} Encombrement (crowding)~; skew porté par la demande de couverture plutôt que par l'information.
\sfield{Horizon, capacité, infrastructure} Quelques semaines à quelques mois~; surfaces de volatilité des options par action (livre~5).
\sfield{Backtest honnêtement} Volatilités implicites issues des cotations de clôture avant la transaction sur l'action, avec leur bid--ask~; le calendrier des événements connu à l'avance.
\sfield{Sources} Xing, Zhang et Zhao (2010)~: 10.9~\% par an, ajustés du risque, entre les portefeuilles aux smirks extrêmes.
\end{strategyfile}

\begin{strategyfile}{Signal d'écart de volatilité implicite}\label{strat:s1:options-implied-signals-for-stocks:spread}
\sfield{Qui paie, et pourquoi} Comme pour le skew, des deux côtés~: des calls chers signalent de bonnes nouvelles, des puts chers de mauvaises.
\sfield{Instruments et plateformes} Actions dont les options sont liquides, à prix d'exercice appariés.
\sfield{Signal} Volatilité implicite du call moins celle du put, pondérée par les positions ouvertes sur l'ensemble des prix d'exercice.
\sfield{Taille et exécution} Acheter les écarts élevés, vendre les faibles~; hebdomadaire ou autour des événements.
\sfield{Coûts} Comme pour le skew.
\sfield{Comment elle s'éteint} Elle s'est affaiblie sur l'échantillon de Cremers et Weinbaum.
\sfield{Horizon, capacité, infrastructure} Quelques jours à quelques semaines.
\sfield{Backtest honnêtement} Primes d'exercice anticipé retirées des options américaines avant le calcul de l'écart.
\sfield{Sources} Cremers et Weinbaum (2010)~: 50 points de base par semaine~; Bali et Hovakimian (2009).
\end{strategyfile}

\begin{strategyfile}{Signal de volume d'options}\label{strat:s1:options-implied-signals-for-stocks:volume}
\sfield{Qui paie, et pourquoi} Traders informés dont le volume d'options monte avant les nouvelles, davantage avant les mauvaises.
\sfield{Instruments et plateformes} Actions pour lesquelles des données de volume d'options sont disponibles.
\sfield{Signal} Le \omterm{def:s1:options-implied-signals-for-stocks:volume}{ratio de volume options sur actions} (non signé) ou un \omterm{def:s1:options-implied-signals-for-stocks:volume}{ratio put--call} des achats à l'ouverture (signé).
\sfield{Taille et exécution} Vendre les ratios élevés~; avec des données signées, acheter les \omterm{def:s1:options-implied-signals-for-stocks:volume}{ratios put--call} faibles.
\sfield{Coûts} Comme ci-dessus.
\sfield{Comment elle s'éteint} Volume lié à la couverture et aux produits structurés, qui noie le flux informé.
\sfield{Horizon, capacité, infrastructure} De quelques jours à une semaine~; volume par série d'options.
\sfield{Backtest honnêtement} Volume publié le jour même, sans révision~; volume signé seulement là où les données sont réellement signées.
\sfield{Sources} Pan et Poteshman (2006)~: plus de 1~\% par semaine~; Johnson et So (2012)~: 0.34~\% par semaine entre les déciles extrêmes.
\end{strategyfile}

\begin{strategyfile}{Signal de volatilité implicite moins réalisée}\label{strat:s1:options-implied-signals-for-stocks:ivrv}
\sfield{Qui paie, et pourquoi} Investisseurs qui paient pour une assurance contre la volatilité, et prime de risque que cette assurance comporte.
\sfield{Instruments et plateformes} Actions et leurs options.
\sfield{Signal} Volatilité réalisée moins implicite (convention de signe de Bali et Hovakimian).
\sfield{Taille et exécution} Tris en coupe transversale, mensuels.
\sfield{Coûts} Faible rotation.
\sfield{Comment elle s'éteint} Comme prime de risque, elle s'éteint lors des pics de volatilité.
\sfield{Horizon, capacité, infrastructure} Quelques mois~; surfaces de volatilité.
\sfield{Backtest honnêtement} Volatilité réalisée issue de données antérieures à la cotation implicite.
\sfield{Sources} Bali et Hovakimian (2009)~; contrôle du chapitre (IC nul là où rien n'est injecté).
\end{strategyfile}

\section{Tutoriel~: quelqu'un achète une protection}\label{tut:s1:options-implied-signals-for-stocks:tut}

\begin{tutorial}
\textbf{Objectif.} Ajouter une couche d'options avec des traders informés au marché synthétique, calculer les signaux d'options et les mesurer par horizon et autour des
événements. \textbf{État final~:} le tableau et la \cref{fig:s1:options-implied-signals-for-stocks:ic}.
\begin{tutsteps}
\item \textbf{La couche}~: volatilités implicites et volumes, mus par la demande informée avant les annonces.
\omcode{code/firm/optsignal/firm_optsignal.py}{37}{58}{Demande d'options informée et signaux qu'elle laisse.}\label{lst:s1:options-implied-signals-for-stocks:layer}
\item \textbf{Lancer} \texttt{ic(name, h)} et \texttt{ic(name, h, True)} pour les quatre signaux et les quatre horizons, \texttt{book(name)} et
\texttt{book(name, True)}, et \texttt{fig\_options.py}.
\end{tutsteps}
\textbf{Pour aller plus loin.} Faire en sorte que l'essentiel du volume d'options soit de la couverture et du bruit, la part du flux informé restant faible~; lisser les signaux pour réduire la rotation~; ajouter un
coût de la vente à découvert et laisser les traders informés choisir entre l'action et les options.
\end{tutorial}

\section{Construction~: la couche d'options}\label{bld:s1:options-implied-signals-for-stocks:optsignal}

\begin{build}
\bfield{Rôle} Un marché d'options synthétique rattaché au panel d'actions, et les signaux issus des options qui y sont construits.
\bfield{Interface} \texttt{OptionConfig}, \texttt{simulate\_options(flag, surprise, listed, rv, cfg, rng)}, \texttt{signals(opt, rv)}.
\bfield{Règles} Demande informée seulement avant les événements et dans le sens de leur surprise~; tout le reste est du bruit~; signaux calculés à partir des valeurs de clôture du jour.
\bfield{Tests d'acceptation} \texttt{code/firm/optsignal/tests/}~: les bonnes options bougent, de la bonne ampleur, seulement dans la fenêtre qui précède l'événement~; le volume monte
davantage pour les mauvaises nouvelles~; la volatilité implicite moins la réalisée égale la prime en l'absence de bruit.
\bfield{Pour aller plus loin} Une surface de volatilité complète par action (livre~5)~; flux de couverture et de produits structurés~; volume signé.
\end{build}

\begin{omsources}
\item Y.~Xing, X.~Zhang et R.~Zhao, ``What does the individual option volatility smirk tell us about future equity returns?'', \emph{Journal of Financial and
Quantitative Analysis} 45(3), 2010.
\item M.~Cremers et D.~Weinbaum, ``Deviations from put-call parity and stock return predictability'', \emph{Journal of Financial and Quantitative Analysis}
45(2), 2010.
\item J.~Pan et A.~M.~Poteshman, ``The information in option volume for future stock prices'', \emph{Review of Financial Studies} 19(3), 2006.
\item T.~L.~Johnson et E.~C.~So, ``The option to stock volume ratio and future returns'', \emph{Journal of Financial Economics} 106(2), 2012.
\item T.~G.~Bali et A.~Hovakimian, ``Volatility spreads and expected stock returns'', \emph{Management Science} 55(11), 2009.
\end{omsources}

\section{Exercices}

\begin{exercise}[$\star$]\label{exo:s1:options-implied-signals-for-stocks:1}
Dans la couche synthétique, évaluer de combien une surprise de $-2$ écarts-types fait monter la volatilité implicite du put dans les jours qui précèdent l'annonce~; et
celle du call pour une surprise de $+2$.
\end{exercise}

\begin{exercise}[$\star$]\label{exo:s1:options-implied-signals-for-stocks:2}
Calculer l'équivalent annuel, sans capitalisation, de l'écart de 0.34~\% par semaine de Johnson et So~; et celui des 50 points de base par semaine de Cremers et Weinbaum.
\end{exercise}

\begin{exercise}[$\star$]\label{exo:s1:options-implied-signals-for-stocks:3}
Expliquer pourquoi le signe du skew, comme signal, est négatif, et celui de l'écart call--put positif~?
\end{exercise}

\begin{exercise}[$\star\star$]\label{exo:s1:options-implied-signals-for-stocks:4}
Expliquer pourquoi les traders informés peuvent préférer les options pour les mauvaises nouvelles plutôt que pour les bonnes~?
\end{exercise}

\begin{exercise}[$\star\star$]\label{exo:s1:options-implied-signals-for-stocks:5}
Les IC sont environ six fois plus élevés parmi les actions qui annoncent dans les dix jours. Expliquer pourquoi, et ce que cela suggère pour le portefeuille~?
\end{exercise}

\begin{exercise}[$\star\star$]\label{exo:s1:options-implied-signals-for-stocks:6}
Expliquer pourquoi le portefeuille hebdomadaire sur l'ensemble des actions perd après coûts alors que son ratio de Sharpe brut est supérieur à 3~?
\end{exercise}

\begin{exercise}[$\star\star\star$]\label{exo:s1:options-implied-signals-for-stocks:7}
\emph{Programmation.} Lancer \texttt{book('os', True)} et le comparer aux portefeuilles fondés sur les prix. Expliquer pourquoi le volume est ici le signal le plus faible~?
\end{exercise}

\begin{exercise}[$\star\star\star$]\label{exo:s1:options-implied-signals-for-stocks:8}
\emph{Trouver la faille.} «~Le portefeuille fondé sur l'écart de volatilité rapporte 32~\% par an après coûts sur le marché synthétique~; on doit s'attendre à en retrouver une large part en pratique.~»
\end{exercise}

\section{Problème~: quelqu'un achète une protection}

\begin{problem}[{problème du week-end --- information dans le marché des options}]\label{pb:s1:options-implied-signals-for-stocks:1}
La couche d'options du chapitre et les archives publiques.

\medskip\noindent\textbf{Partie I --- Les signaux.}
\begin{enumerate}
\item Définir le \omterm{def:s1:options-implied-signals-for-stocks:informed}{trading d'options informé}, l'\omterm{def:s1:options-implied-signals-for-stocks:spread}{écart de volatilité implicite}, le \omterm{def:s1:options-implied-signals-for-stocks:volume}{ratio put--call} et le \omterm{def:s1:options-implied-signals-for-stocks:volume}{ratio de volume options sur actions}.
\item Que Xing, Zhang et Zhao ont-ils constaté au sujet du smirk~?
\item Que Cremers et Weinbaum, puis Bali et Hovakimian, ont-ils constaté au sujet des spreads~?
\item Que Pan et Poteshman, puis Johnson et So, ont-ils constaté au sujet du volume~?
\end{enumerate}

\medskip\noindent\textbf{Partie II --- La couche.}
\begin{enumerate}[resume]
\item Décrire la couche d'options synthétiques et ses traders informés.
\item Pourquoi les mauvaises nouvelles sont-elles davantage pondérées~?
\item Pourquoi la différence entre volatilité implicite et volatilité réalisée sert-elle ici de contrôle~?
\item Comment le calendrier des événements est-il utilisé~?
\end{enumerate}

\medskip\noindent\textbf{Partie III --- Les mesures.}
\begin{enumerate}[resume]
\item Donner les IC par horizon pour l'ensemble des actions.
\item Donner-les pour les actions annonçant dans les dix jours.
\item Pourquoi les IC culminent-ils à vingt jours~?
\item Donner les résultats des livres, pour l'ensemble des actions et autour des événements.
\end{enumerate}

\medskip\noindent\textbf{Partie IV --- Le verdict.}
\begin{enumerate}[resume]
\item Énoncer le \emph{résultat nommé}~: l'IC du skew et des signaux d'écart de volatilité, et leur décroissance avec l'horizon.
\item Pourquoi les rendements synthétiques sont-ils très supérieurs à ceux de la littérature~?
\item Que modifieriez-vous pour rendre la couche réaliste~?
\item Comment réduiriez-vous la rotation du livre portant sur l'ensemble des actions~?
\item De quelles données auriez-vous besoin pour négocier ces stratégies~?
\item Quel signal s'est affaibli dans les publications~?
\item Comment ces signaux se rapportent-ils aux stratégies sur résultats du chapitre~7~?
\item En une phrase~: pourquoi regarder les options pour négocier des actions~?
\end{enumerate}
\end{problem}

\section{Questions d'entretien}

\begin{interviewq}[$\star$ \iqroles{chercheur, trader}]\label{iq:s1:options-implied-signals-for-stocks:1}
Pourquoi les prix des options pourraient-ils prédire les rendements des actions~?
\end{interviewq}

\begin{interviewq}[$\star\star$ \iqroles{chercheur}]\label{iq:s1:options-implied-signals-for-stocks:2}
Comment calculeriez-vous le skew de volatilité d'une action pour un signal en coupe transversale~?
\end{interviewq}

\begin{interviewq}[$\star\star$ \iqroles{chercheur}]\label{iq:s1:options-implied-signals-for-stocks:3}
Quelle est la différence, comme signaux, entre l'\omterm{def:s1:options-implied-signals-for-stocks:spread}{écart de volatilité implicite} call--put et la différence entre volatilité implicite et volatilité réalisée~?
\end{interviewq}

\begin{interviewq}[$\star\star$ \iqroles{trader}]\label{iq:s1:options-implied-signals-for-stocks:4}
Le volume de puts d'une action triple avant les résultats. Que faites-vous~?
\end{interviewq}

\begin{interviewq}[$\star\star$ \iqroles{risque}]\label{iq:s1:options-implied-signals-for-stocks:5}
Quels sont les pièges du backtest de signaux issus des options~?
\end{interviewq}

\begin{interviewq}[$\star\star\star$ \iqroles{chercheur}]\label{iq:s1:options-implied-signals-for-stocks:6}
Dans un modèle où les traders informés choisissent entre vendre l'action à découvert à un coût $c$ et acheter des puts avec un effet de levier $L$, déterminer quand ils préfèrent les puts, et
ce que cela implique pour l'information contenue dans le volume de puts.
\end{interviewq}
